% ==============================================================================
\section{Trabajo a futuro}
\label{sec:trabajo_futuro}
% ==============================================================================

El trabajo pendiente se ordena según su efecto sobre la seguridad del usuario y
según lo que se aprendió en las mediciones de esta etapa.

\textbf{Correcciones identificadas en el software.} Varias deficiencias quedaron
localizadas con precisión y su corrección es acotada:

\begin{itemize}
    \item Exigir que la llegada por percepción se cumpla durante varios cuadros
    seguidos, o combinarla con la distancia en el mapa, para que un error monocular
    aislado no termine la sesión.
    \item Agrupar el tiempo de espera de las alertas por tipo y sector, para que una
    precaución lateral no silencie una frontal, y unificar la ortografía de las frases.
    \item Permitir que los rayos precisos reduzcan el log-odds de las celdas
    consolidadas y que la capa de peligros pierda impactos con el tiempo, de modo que
    los obstáculos que desaparecen se borren del mapa (RF-17).
    \item Traducir al español las etiquetas de los \emph{landmarks} en los mensajes de
    llegada y en la descripción del entorno.
    \item Ampliar las frases de referencia del PLN con las expresiones que fallaron en la
    evaluación y revisar la relación entre ``ayuda'' y ``emergencia''.
    \item Hacer que la petición de ruta alternativa obligue al agente a evitar el
    camino anterior.
\end{itemize}

\textbf{Integración del hardware.} Montar los componentes en el armazón diseñado,
conectar la IMU BNO055 y los sensores MB1232 por I2C, medir con regla la geometría
real y ajustarla en \texttt{montaje.py}, ejecutar la calibración de la IMU con los
lentes puestos y calibrar la longitud de paso de cada usuario. Con
los sensores montados se debe medir el ruido real de la odometría, ajustar el
simulador y reentrenar el agente de navegación. También hay que medir los tiempos de
cada módulo en la Raspberry Pi y, si no alcanzan los requerimientos, probar los
modelos de visión en sus versiones más pequeñas o la arquitectura con una computadora
externa que se planteó en TT1.

\textbf{Percepción con datos propios.} Etiquetar imágenes de los pasillos, puertas y
mamparas de la escuela para reentrenar el detector con el script preparado y, con el
mismo conjunto, medir la precisión que exige RNF-06. Separar la banqueta del arroyo
vehicular en la segmentación, para que el sistema no considere transitable la calle.

\textbf{Localización.} Agregar cierre de lazo o una prueba de compatibilidad conjunta
en la asociación de datos para evitar los \emph{landmarks} duplicados y la
divergencia de orientación observada en recorridos largos.

\textbf{Voz en condiciones reales.} Repetir la evaluación de reconocimiento con
personas reales, con el micrófono del prototipo y en lugares con ruido de calle, y
probar un modelo de Vosk más grande si la Raspberry Pi lo permite. Implementar el
aviso a un contacto de emergencia.

\textbf{Pruebas con usuarios.} La validación final del sistema requiere pruebas con
personas con discapacidad visual, en escenarios controlados, acompañadas y con el
bastón blanco en uso. Antes de ello conviene probar el sistema con personas con los
ojos vendados en un entorno conocido, para ajustar el ritmo y la redacción de las
indicaciones.
